\documentclass[12pt,a4paper]{article}
\usepackage[margin=1in]{geometry}
\usepackage[T1]{fontenc}
\usepackage{lmodern}
\usepackage{setspace}
\usepackage{graphicx}
\usepackage{float}
\usepackage{hyperref}
\usepackage{longtable}
\usepackage{array}
\usepackage{listings}
\usepackage{xcolor}
\usepackage{titlesec}

\hypersetup{
    colorlinks=true,
    linkcolor=blue,
    urlcolor=blue
}

\onehalfspacing
\setlength{\parskip}{0.5em}
\setlength{\parindent}{0pt}

\titleformat{\section}{\large\bfseries}{\thesection.}{0.5em}{}
\titleformat{\subsection}{\normalsize\bfseries}{\thesubsection}{0.5em}{}

\lstdefinestyle{cppstyle}{
    language=C++,
    basicstyle=\ttfamily\small,
    keywordstyle=\color{blue},
    commentstyle=\color{gray},
    stringstyle=\color{teal},
    numbers=left,
    numberstyle=\tiny\color{gray},
    stepnumber=1,
    numbersep=8pt,
    frame=single,
    breaklines=true,
    showstringspaces=false,
    tabsize=4
}

\begin{document}

%--------------------------
% Title Page
%--------------------------
\begin{titlepage}
    \centering
    \vspace*{2cm}
    {\Huge \textbf{Chess Engine with SFML GUI}}\\[0.8cm]
    {\Large Competitive Programming Lab Mini Project Report}\\[2cm]

    \begin{tabular}{>{\bfseries}p{5cm} p{8cm}}
        Student Name: & \underline{\hspace{8cm}} \\
        Roll Number: & \underline{\hspace{8cm}} \\
        Program/Semester: & \underline{\hspace{8cm}} \\
        Course: & Competitive Programming Lab \\
        Course Code: & 24B15CS224 \\
        Faculty: & \underline{\hspace{8cm}} \\
        Session: & 2025--2026 (Even Semester) \\
        Submission Date: & \underline{\hspace{8cm}} \\
    \end{tabular}

    \vfill
    {\large Department of Computer Science}\\
    {\large \underline{\hspace{8cm}}}
\end{titlepage}

\tableofcontents
\newpage

%--------------------------
% Summary
%--------------------------
\section{Summary}
\subsection{Brief Overview of the Project}
This project implements a desktop chess application in C++ using SFML for rendering and interaction. The system supports legal chess moves, turn-based play (human vs AI), check/checkmate/stalemate detection, and a minimax-based AI enhanced with alpha-beta pruning. It also includes benchmarking output and visual feedback features for better user experience.

\subsection{Objectives}
\begin{itemize}
    \item Build a complete playable chess system with graphical interface.
    \item Implement rule-based move legality and game state detection.
    \item Integrate AI using minimax and alpha-beta pruning.
    \item Compare minimax and alpha-beta using performance metrics.
    \item Improve user interaction using highlights and status panels.
\end{itemize}

%--------------------------
% Introduction
%--------------------------
\section{Introduction}
\subsection{Background and Context}
Chess is a classical two-player strategy game often used for teaching algorithmic thinking, state-space search, and optimization. This project applies core competitive programming concepts (data structures, search trees, pruning, complexity analysis) in a practical mini-project.

\subsection{Problem Statement}
Design and implement an interactive chess engine that:
\begin{itemize}
    \item validates player moves correctly,
    \item plays as an intelligent AI opponent,
    \item provides clear visual feedback and game status information.
\end{itemize}

\subsection{Project Objectives}
\begin{itemize}
    \item Represent board state efficiently.
    \item Generate legal moves and enforce rules.
    \item Evaluate and choose AI moves using minimax.
    \item Optimize search with alpha-beta pruning.
    \item Present game state and analytics in a clean SFML interface.
\end{itemize}

\subsection{Scope of the Project}
Included:
\begin{itemize}
    \item Board rendering, piece rendering, move interaction.
    \item Core move legality for all piece types.
    \item Check/checkmate/stalemate detection.
    \item Minimax + alpha-beta + benchmark output.
    \item Move quality classification (Best/Good/Inaccuracy/Mistake/Blunder).
\end{itemize}

Not included:
\begin{itemize}
    \item Online multiplayer.
    \item Opening book/endgame tablebases.
    \item Advanced heuristics beyond current evaluation function.
\end{itemize}

%--------------------------
% System Requirements
%--------------------------
\section{System Requirements}
\subsection{Functional Requirements}
\begin{itemize}
    \item Open and render a chessboard window.
    \item Load piece textures and move sound assets.
    \item Allow piece selection and valid destination moves.
    \item Reject invalid moves with feedback.
    \item Alternate turns between White (player) and Black (AI).
    \item Detect check, checkmate, and stalemate.
    \item Compute AI moves using minimax/alpha-beta.
    \item Display benchmark metrics: nodes explored and time.
    \item Display move quality and game status in side panel.
\end{itemize}

%--------------------------
% Design and Implementation
%--------------------------
\section{Design and Implementation}
\subsection{UML Diagrams}
\textbf{Class/Module-Level Design (simplified):}
\begin{itemize}
    \item \texttt{board.h / board.cpp}: board type, initial board, print board.
    \item \texttt{moves.h / moves.cpp}: move legality, check detection, move generation, make/undo.
    \item \texttt{engine.h / engine.cpp}: evaluation, minimax, alpha-beta, best move, benchmarking counters.
    \item \texttt{main.cpp}: SFML window, input handling, rendering, side panel, integration.
\end{itemize}

\textbf{Interaction Flow:}
\begin{enumerate}
    \item User selects source and destination squares.
    \item System validates move using rule engine.
    \item If valid, move executes and quality is evaluated.
    \item AI computes response using minimax/alpha-beta.
    \item UI updates highlights, status, benchmark output.
\end{enumerate}

\subsection{Algorithms Implementation}
\textbf{Implemented algorithms:}
\begin{itemize}
    \item Legal move generation (piece-wise move rules + king safety).
    \item Minimax search (plain and alpha-beta variants).
    \item Move ordering via heuristic score and \texttt{lower\_bound}.
    \item Board evaluation based on material and positional bonuses.
\end{itemize}

\subsection{Algorithms and Their Complexities}
\begin{longtable}{|p{5cm}|p{4cm}|p{5cm}|}
\hline
\textbf{Algorithm} & \textbf{Time Complexity} & \textbf{Notes} \\
\hline
Move generation (single position) & $O(B \times 64)$ approx. & Tries destinations for each piece, filtered by legality \\
\hline
Minimax (depth d) & $O(b^d)$ & $b$ = branching factor \\
\hline
Alpha-beta (depth d) & Best better than minimax, worst $O(b^d)$ & Prunes subtrees; practical speedup \\
\hline
Move ordering insertion (\texttt{lower\_bound}) & $O(\log n)$ search + $O(n)$ insert & For sorted scored move list \\
\hline
Check detection & $O(64 \times M)$ & Tests opponent pseudo-legal attacks on king square \\
\hline
\end{longtable}

\subsection{Code Snippets}
\textbf{Minimax with Alpha-Beta (excerpt):}
\begin{lstlisting}[style=cppstyle]
int minimax_ab(Board& board, int depth, int alpha, int beta, bool isMaximizing) {
    ++nodesExplored;
    if (depth == 0) return evaluateBoard(board);

    vector<Move> moves = getOrderedMoves(board, isMaximizing);
    if (moves.empty()) return evaluateBoard(board);

    if (isMaximizing) {
        int bestScore = INT_MIN;
        for (const Move& move : moves) {
            char captured = makeMove(board, move);
            int score = minimax_ab(board, depth - 1, alpha, beta, false);
            undoMove(board, move, captured);
            bestScore = max(bestScore, score);
            alpha = max(alpha, bestScore);
            if (beta <= alpha) break;
        }
        return bestScore;
    }
    int bestScore = INT_MAX;
    for (const Move& move : moves) {
        char captured = makeMove(board, move);
        int score = minimax_ab(board, depth - 1, alpha, beta, true);
        undoMove(board, move, captured);
        bestScore = min(bestScore, score);
        beta = min(beta, bestScore);
        if (beta <= alpha) break;
    }
    return bestScore;
}
\end{lstlisting}

\textbf{Move quality classification (excerpt):}
\begin{lstlisting}[style=cppstyle]
int centipawnLoss = max(0, (bestMoveEval - playedMoveEval) * 10);
string moveQuality = "Blunder";
if (centipawnLoss <= 20) moveQuality = "Best Move";
else if (centipawnLoss <= 50) moveQuality = "Good Move";
else if (centipawnLoss <= 100) moveQuality = "Inaccuracy";
else if (centipawnLoss <= 300) moveQuality = "Mistake";
\end{lstlisting}

\subsection{Output Screenshots with Description}
\begin{figure}[H]
    \centering
    \includegraphics[width=0.9\textwidth]{figures/gameplay_main.png}
    \caption{Main gameplay screen with board, pieces, and highlights.}
\end{figure}

\begin{figure}[H]
    \centering
    \includegraphics[width=0.9\textwidth]{figures/side_panel_status.png}
    \caption{Side panel showing game status, move quality, and evaluations.}
\end{figure}

\begin{figure}[H]
    \centering
    \includegraphics[width=0.9\textwidth]{figures/benchmark_output.png}
    \caption{Console output comparing minimax and alpha-beta benchmark metrics.}
\end{figure}

\textit{Note: Replace image paths with actual screenshots saved in a \texttt{figures/} folder.}

%--------------------------
% Conclusion
%--------------------------
\section{Conclusion}
\subsection{Summary \& Achievement of Objectives}
The project successfully delivers a working SFML chess application with legal move validation, AI move generation, and rich gameplay feedback. Core objectives such as implementing minimax with alpha-beta pruning, benchmark comparison, and UI improvements were achieved.

\subsection{Future Work and Recommendations}
\begin{itemize}
    \item Add advanced chess rules (castling, en passant, promotion logic enhancements if needed).
    \item Improve evaluation with piece-square tables and king safety terms.
    \item Add transposition table and iterative deepening for stronger AI.
    \item Add PGN/FEN support for reproducible test positions.
    \item Add automated test suite for move generation and checkmate scenarios.
\end{itemize}

%--------------------------
% References
%--------------------------
\section{References}
\begin{enumerate}
    \item Steven Halim et al., \textit{Competitive Programming 3}, Lulu Independent Publish, 2013.
    \item Antti Laaksonen, \textit{Guide to Competitive Programming}, Springer, 2020.
    \item Antti Laaksonen, \textit{Competitive Programmer's Handbook}, 2017.
    \item SFML Documentation, \url{https://www.sfml-dev.org/documentation/}
    \item C++ Reference, \url{https://en.cppreference.com/}
\end{enumerate}

\end{document}
